% !PRINTED 539
\weberpage{539}

%<<CONT
\weq{16}{\begin{aligned}
\operatorname*{Lim}_{s=1} (4^s + 2 - 3.\,2^s) \overset{\psi}{\Sigma} S &= \frac{24\pi}{\sqrt{m}} \log \Pi \frac{f(\omega)}{\sqrt[4]{2}}, & m &\equiv 3 \pmod 8, \\
\operatorname*{Lim}_{s=1} (4^s + 2 - 3.\,2^s) \overset{\psi}{\Sigma} S &= \frac{8\pi}{\sqrt{m}} \log \Pi \frac{f(\omega)}{\sqrt[4]{2}}, & m &\equiv 7 \pmod 8.
\end{aligned}}
But we have [§ 141, (28)]:
\[
\mathrm{Lim}\,(s-1) \overset{\psi}{\Sigma} S = \frac{\pi h}{\sqrt{m}}, \quad \mathrm{Lim}\,\frac{4^s + 2 - 3.\,2^s}{s-1} = 2 \log 2,
\]
hence:
\weq{17}{\begin{aligned}
\log \Pi \frac{f(\omega)}{\sqrt[3]{2}} &= 0, & m &\equiv 3 \pmod 8, \\
\log \Pi \frac{f(\omega)}{\sqrt{2}} &= 0, & m &\equiv 7 \pmod 8.
\end{aligned}}
The first of these formulae gives the result already derived in another
way; the second gives the new one:
\weq{18}{\tau = \tfrac{1}{2}, \quad m \equiv 7 \pmod 8.}
4. There still remains the determination of $\tau$ in the case where
$m$ is divisible by a higher power of 2. In order to carry out this
determination as well, let
\weq{19}{m = 4m',}
% [sic] The source prints sqrt(m); (19), m'=AC-B^2 and the root of (A,2B,C) require sqrt(m') here.
\weq{20}{\omega' = \frac{-B + i\sqrt{m}}{A}, \quad m' = AC - B^2,}
so that $\omega'$ is the root of the quadratic form $(A, 2B, C)$ of
discriminant $-m$, in which $A$ and $C$ may be assumed to be \gesperrt{odd}. Then
\weq{21}{\omega_1 = 2\,\omega', \quad \omega_2 = \frac{\omega'}{2}}
are roots of the forms
\weq{22}{(A, 4B, 4C), \quad (4A, 4B, C),}
and by § 126, (13) the 24th powers of
\[
f_1(\omega_1) = \frac{\sqrt{2}}{f_2(\omega')}, \quad f_2(\omega_2) = \frac{\sqrt{2}}{f_1(\omega')}
\]
are class invariants of the discriminant $-4m$.

If $\omega'$ runs through a complete system of roots of the discriminant $-m$, then $\omega_1$ and $\omega_2$ together run through a complete
system of roots of the discriminant $-4m$ (§ 123), and we obtain,
%CONT>>
